\documentclass[12pt]{article}
\usepackage{amsmath,amssymb,amsthm}
\usepackage[margin=1in]{geometry}

\newtheorem{theorem}{Theorem}
\newtheorem{lemma}[theorem]{Lemma}
\newtheorem{corollary}[theorem]{Corollary}
\newtheorem{proposition}[theorem]{Proposition}
\theoremstyle{definition}
\newtheorem{definition}[theorem]{Definition}

\title{Project Euler Problem 419: Look and Say Sequence}
\author{Solution}
\date{}

\begin{document}
\maketitle

\section{Problem Statement}
The look-and-say sequence begins 1, 11, 21, 1211, 111221, 312211, \ldots

Define $A(n)$, $B(n)$, $C(n)$ as the count of ones, twos, and threes in the $n$-th term. Find $A(n), B(n), C(n)$ for $n = 10^{12}$, modulo $2^{30}$.

\section{Mathematical Analysis}

\subsection{Conway's Cosmological Theorem}
\begin{theorem}[Conway, 1986]
Every look-and-say string eventually splits into a compound of exactly 92 ``elements'' (atoms), each evolving independently under the look-and-say operation.
\end{theorem}

\subsection{Transition Matrix}
Let $M$ be the $92 \times 92$ transition matrix where $M_{ij}$ counts how many copies of element $j$ appear in the decay of element $i$. Then the element composition vector evolves as:
\[
\mathbf{v}(n) = M^{n-k} \cdot \mathbf{v}(k)
\]

\subsection{Digit Counting}
Define weight vectors $\mathbf{w}_d$ ($d = 1, 2, 3$) where $(\mathbf{w}_d)_i$ counts occurrences of digit $d$ in element $i$. Then:
\[
A(n) = \mathbf{w}_1^T \cdot \mathbf{v}(n), \quad
B(n) = \mathbf{w}_2^T \cdot \mathbf{v}(n), \quad
C(n) = \mathbf{w}_3^T \cdot \mathbf{v}(n)
\]

\subsection{Matrix Exponentiation}
Since $n = 10^{12}$, we compute $M^n \pmod{2^{30}}$ via repeated squaring:
\begin{itemize}
    \item $\log_2(10^{12}) \approx 40$ squarings
    \item Each squaring multiplies two $92 \times 92$ matrices
\end{itemize}

\section{Algorithm}
\begin{enumerate}
    \item Define Conway's 92 elements and the transition matrix $M$.
    \item Compute initial composition vector $\mathbf{v}(k)$ for small $k$.
    \item Compute $M^{n-k} \pmod{2^{30}}$ via matrix exponentiation.
    \item Extract $A(n), B(n), C(n)$ using digit-weight vectors.
\end{enumerate}

\section{Complexity}
\begin{itemize}
    \item \textbf{Time:} $O(92^3 \cdot \log n)$
    \item \textbf{Space:} $O(92^2)$
\end{itemize}

\section{Answer}

\[
\boxed{998567458,1046245404,43363922}
\]

\end{document}
